% =============================================================================
% Chapter: RAG & Retrieval-based LLMs
% =============================================================================

\chapter{RAG \& Retrieval-based LLMs}

\begin{formulabox}[Key Concepts]
\textbf{RAG (Retrieval-Augmented Generation):}

Συνδυασμός retrieval από knowledge base + generation από LLM.

\textbf{Basic Pipeline:}
\begin{enumerate}
    \item \textbf{Query}: User question
    \item \textbf{Retrieve}: Find relevant documents
    \item \textbf{Augment}: Add context to prompt
    \item \textbf{Generate}: LLM produces answer
\end{enumerate}

\textbf{Key Metrics:}
\begin{itemize}
    \item Recall@K: Fraction of relevant docs in top-K
    \item Precision@K: Fraction of top-K that are relevant
    \item MRR (Mean Reciprocal Rank): $\frac{1}{|Q|}\sum_{i=1}^{|Q|} \frac{1}{\text{rank}_i}$
\end{itemize}
\end{formulabox}

% =============================================================================
\section{RAG Architecture}

\begin{center}
\begin{tikzpicture}[
    block/.style={draw=primary, fill=box-blue, rounded corners=5pt, minimum width=2.5cm, minimum height=1cm, font=\small},
    arrow/.style={->, thick}
]
    % Query
    \node[block, fill=secondary!30] (query) at (0,0) {Query};

    % Retriever
    \node[block] (retriever) at (4,0) {Retriever};

    % Knowledge Base
    \node[block, fill=success!30, minimum height=2cm, align=center] (kb) at (4,-2.5) {Knowledge\\Base};

    % Context
    \node[block] (context) at (8,0) {Context};

    % LLM
    \node[block, fill=accent!30] (llm) at (12,0) {LLM};

    % Answer
    \node[block, fill=success!30] (answer) at (12,-2) {Answer};

    % Connections
    \draw[arrow] (query) -- (retriever);
    \draw[arrow, secondary] (kb) -- (retriever);
    \draw[arrow] (retriever) -- (context);
    \draw[arrow] (query.east) to[bend left=20] (context.north);
    \draw[arrow] (context) -- (llm);
    \draw[arrow] (llm) -- (answer);

    % Labels
    \node[font=\footnotesize, anchor=north] at (4,-0.6) {Search};
    \node[font=\footnotesize, anchor=south] at (10,0.1) {Augmented Prompt};
\end{tikzpicture}
\end{center}

% =============================================================================
\section{Retrieval Methods}

\subsection{Sparse Retrieval (Traditional)}

\begin{definitionbox}[TF-IDF / BM25]
Keyword-based retrieval using term frequencies.

\textbf{TF-IDF:}
\[
\text{TF-IDF}(t, d) = TF(t, d) \times \log\left(\frac{N}{DF(t)}\right)
\]

\textbf{BM25:} Improved TF-IDF with document length normalization.

\textbf{Pros:} Fast, interpretable, no training required\\
\textbf{Cons:} No semantic understanding (keyword mismatch problem)
\end{definitionbox}

\subsection{Dense Retrieval (Neural)}

\begin{definitionbox}[Dense Embeddings]
Documents and queries encoded as vectors in semantic space.

\textbf{Process:}
\begin{enumerate}
    \item Encode documents → $\mathbf{d}_i = \text{Encoder}(d_i)$
    \item Encode query → $\mathbf{q} = \text{Encoder}(q)$
    \item Score = $\cos(\mathbf{q}, \mathbf{d}_i)$ or $\mathbf{q} \cdot \mathbf{d}_i$
\end{enumerate}

\textbf{Pros:} Captures semantic similarity\\
\textbf{Cons:} Requires training, less interpretable
\end{definitionbox}

\begin{center}
\begin{tikzpicture}
    % Sparse
    \node[draw=secondary, fill=box-blue, rounded corners=8pt, minimum width=5cm, minimum height=3cm] (sparse) at (0,0) {};
    \node[anchor=north, font=\bfseries] at (sparse.north) {Sparse (TF-IDF/BM25)};
    \node[anchor=center, align=center, font=\small] at (sparse.center) {
        Keyword matching\\
        Exact terms\\
        Bag-of-words
    };

    % vs
    \node at (3.5,0) {\textbf{vs}};

    % Dense
    \node[draw=success, fill=box-green, rounded corners=8pt, minimum width=5cm, minimum height=3cm] (dense) at (7,0) {};
    \node[anchor=north, font=\bfseries] at (dense.north) {Dense (Embeddings)};
    \node[anchor=center, align=center, font=\small] at (dense.center) {
        Semantic matching\\
        Vector similarity\\
        Context-aware
    };
\end{tikzpicture}
\end{center}

% =============================================================================
\section{Vector Databases}

\begin{infobox}[Vector Similarity Search]
For efficient retrieval from millions of embeddings:

\textbf{Exact Search:}
\begin{itemize}
    \item Brute-force comparison with all vectors
    \item $O(n)$ complexity
\end{itemize}

\textbf{Approximate Nearest Neighbor (ANN):}
\begin{itemize}
    \item HNSW (Hierarchical Navigable Small World)
    \item IVF (Inverted File Index)
    \item Product Quantization (PQ)
\end{itemize}

Trade-off: Speed vs Accuracy
\end{infobox}

% =============================================================================
\section{Chunking Strategies}

\begin{definitionbox}[Document Chunking]
Splitting documents into smaller pieces for retrieval.

\textbf{Strategies:}
\begin{itemize}
    \item \textbf{Fixed-size:} Every N tokens/characters
    \item \textbf{Sentence-based:} Split at sentence boundaries
    \item \textbf{Paragraph-based:} Natural paragraph breaks
    \item \textbf{Semantic:} Use embeddings to find natural breaks
    \item \textbf{Overlapping:} Chunks with shared context
\end{itemize}

\textbf{Considerations:}
\begin{itemize}
    \item Too small: Lost context
    \item Too large: Irrelevant content included
    \item Typical: 256-512 tokens with 20\% overlap
\end{itemize}
\end{definitionbox}

% =============================================================================
\section{RAG Evaluation}

\begin{center}
\begin{tikzpicture}
    \matrix (m) [
        matrix of nodes,
        nodes={
            draw=primary,
            minimum width=6cm,
            minimum height=1.2cm,
            anchor=center,
            font=\small
        },
        column 1/.style={nodes={fill=ntua-blue!20, font=\bfseries}},
        row 1/.style={nodes={fill=secondary!30, font=\bfseries}},
        column sep=-\pgflinewidth,
        row sep=-\pgflinewidth
    ] {
        Metric & Formula & Use Case \\
        Recall@K & $\frac{|\text{relevant} \cap \text{top-K}|}{|\text{relevant}|}$ & Coverage \\
        Precision@K & $\frac{|\text{relevant} \cap \text{top-K}|}{K}$ & Accuracy \\
        MRR & $\frac{1}{|Q|}\sum\frac{1}{\text{rank}_i}$ & Ranking quality \\
        NDCG & Discounted cumulative gain & Graded relevance \\
    };
\end{tikzpicture}
\end{center}

% =============================================================================
\section{Common Challenges}

\begin{warningbox}[RAG Challenges]
\begin{enumerate}
    \item \textbf{Retrieval Failure:}
    \begin{itemize}
        \item Relevant documents not retrieved
        \item Solution: Better embeddings, hybrid retrieval
    \end{itemize}

    \item \textbf{Context Window Limit:}
    \begin{itemize}
        \item Too much retrieved content
        \item Solution: Re-ranking, summarization
    \end{itemize}

    \item \textbf{Hallucination:}
    \begin{itemize}
        \item LLM ignores or contradicts retrieved context
        \item Solution: Better prompting, fact verification
    \end{itemize}

    \item \textbf{Latency:}
    \begin{itemize}
        \item Retrieval adds overhead
        \item Solution: Caching, ANN indexes
    \end{itemize}
\end{enumerate}
\end{warningbox}

% =============================================================================
\section{Advanced RAG Techniques}

\begin{infobox}[RAG Improvements]
\textbf{Hybrid Retrieval:}
Combine sparse (BM25) + dense retrieval, then merge results.

\textbf{Re-ranking:}
Use cross-encoder to re-score top-K candidates.

\textbf{Query Expansion:}
Use LLM to generate alternative queries.

\textbf{Multi-hop RAG:}
Iterative retrieval for complex questions.
\end{infobox}

% =============================================================================
\section{Quick Reference}

\begin{center}
\begin{tikzpicture}
    \node[draw=ntua-blue, fill=ntua-blue!10, rounded corners=5pt, minimum width=14cm, minimum height=6cm] (main) {};

    \node[anchor=north west, font=\bfseries\Large\color{ntua-blue}] at ([xshift=5pt, yshift=-5pt]main.north west) {RAG Summary};

    \node[anchor=north west, align=left, font=\small] at ([xshift=10pt, yshift=-35pt]main.north west) {
        \textbf{RAG Pipeline:} Query → Retrieve → Augment → Generate\\[5pt]
        \textbf{Retrieval Types:}\\
        • Sparse (BM25): Keyword matching, fast, no training\\
        • Dense (Embeddings): Semantic matching, requires training\\[5pt]
        \textbf{Key Design Choices:}\\
        • Chunk size (256-512 tokens typical)\\
        • Embedding model selection\\
        • Number of retrieved documents (K)
    };

    \node[draw=secondary, fill=box-blue, rounded corners=5pt, minimum width=5cm, minimum height=1.5cm]
        at ([xshift=-4cm, yshift=0cm]main.south) {};
    \node[align=center, font=\footnotesize] at ([xshift=-4cm, yshift=0cm]main.south) {
        \textbf{Why RAG?}\\
        Up-to-date knowledge\\
        Reduces hallucination
    };

    \node[draw=accent, fill=box-red, rounded corners=5pt, minimum width=5cm, minimum height=1.5cm]
        at ([xshift=3cm, yshift=0cm]main.south) {};
    \node[align=center, font=\footnotesize] at ([xshift=3cm, yshift=0cm]main.south) {
        \textbf{Hybrid = Best:}\\
        BM25 + Dense retrieval\\
        then Re-rank
    };
\end{tikzpicture}
\end{center}
